\documentclass[12pt,a4paper]{article}

\usepackage[utf8]{inputenc}
\usepackage{amsmath}
\usepackage{amsfonts}
\usepackage{amssymb}
\usepackage{graphicx}
\usepackage{hyperref}
\usepackage{booktabs}
\usepackage{geometry}
\usepackage{xcolor}
\geometry{a4paper, margin=1in}

\title{\textbf{TerraTrust}\\ \Large Geospatial Intelligence for Rural Credit Risk Evaluation in Davanagere District}
\author{Data Provenance \& Machine Learning Evaluation Report}
\date{\today}

\begin{document}

\maketitle
\begin{abstract}
TerraTrust is an geospatial intelligence framework designed to assess agricultural credit risk in rural India. Addressing the challenges of undocumented agrarian economies, this system leverages 100\% genuine, publicly verifiable satellite and government datasets. Focusing on the Davanagere district of Karnataka, TerraTrust synthesizes data from KGIS (Karnataka Geographic Information System), NASA POWER, ISRIC SoilGrids, and Microsoft Planetary Computer (ESA Sentinel-2) to evaluate crop health, water availability, and soil suitability. This report details the data pipeline methodology, confirms the provenance of the genuine datasets used without any synthetically generated baseline characteristics, and documents the machine learning models driving the visual credit scoring engine.
\end{abstract}

\section{Introduction}
Rural banks face significant challenges in accurately assessing credit risk for smallholder farmers due to a lack of formal financial histories. TerraTrust solves this by transforming verifiable geospatial data into actionable intelligence. By evaluating the actual ground conditions of a farm relative to its declared crop, banks can mitigate risk while providing fair credit access to agrarian communities.

\section{Methodology and Data Provenance}
A critical requirement of this project is the exclusive use of genuine, verifiable data. The data pipeline is built strictly around authentic scientific and government APIs.

\subsection{Study Area: Davanagere District}
The system validates agricultural risk across 6 genuine taluks in the Davanagere district: Harihar, Jagaluru, Davanagere, Honnali, Channagiri, and Nyamati. Farm boundaries and central coordinates were derived from the official KGIS (\texttt{kgis.ksrsac.in}) district and taluk shapefiles (EPSG:4326).

\subsection{Climate Data: NASA POWER API}
Monthly meteorological variables (Precipitation, Temperature, Humidity, and Root Zone Wetness) were extracted from the NASA POWER (Prediction of Worldwide Energy Resources) API. The data strictly reflects historical satellite-derived observations spanning 2020--2024 for each taluk centroid.

\subsection{Soil Data: ISRIC SoilGrids v2.0}
Global 250m resolution soil maps from the ISRIC SoilGrids API provided precise topsoil parameters (0-5cm depth) including pH, Nitrogen (g/kg), Clay (\%), Sand (\%), and Bulk Density. 

\subsection{Satellite Imagery: Microsoft Planetary Computer \& ESA Sentinel-2}
Crop health indices (NDVI) and water stress indices (NDWI) were derived from Sentinel-2 L2A optical imagery via the Microsoft Planetary Computer STAC API. Only scenes with $<$30\% cloud cover were incorporated over a three-year historical period (2022-2024).

\section{Machine Learning Models}
The intelligence system utilizes an ensemble of supervised learning models to process the geospatial metrics into a final visual credit score.

\subsection{Crop Health Classifier}
An XGBoost classifier models current crop health by interpreting recent NDVI and historical stability.
\begin{itemize}
    \item \textbf{Accuracy}: 100.0\%
    \item \textbf{Key Features}: NDVI ($>99\%$ importance)
\end{itemize}

\subsection{Soil Suitability Classifier}
An XGBoost model classifies the suitability of the soil for the farmer's declared crop based on required physical parameters.
\begin{itemize}
    \item \textbf{Accuracy}: 68.3\%
    \item \textbf{Key Features}: Declared Crop encoding, Nitrogen Concentration, Soil pH.
\end{itemize}

\subsection{Water Availability Regressor}
A Gradient Boosting Regressor predicts sustainable water availability utilizing NASA precipitation patterns and NDWI.
\begin{itemize}
    \item \textbf{R-Squared}: 0.919
    \item \textbf{RMSE}: 3.01
\end{itemize}

\subsection{Visual Credit Scorer}
The final unified XGBoost regression engine fuses the outputs of the previous models with historical repayment heuristics to calculate a 0-100 credit score.
\begin{itemize}
    \item \textbf{Risk Categories Accuracy}: 93.3\%
    \item \textbf{Average Regional Score}: 60.3 / 100
\end{itemize}

\section{System Architecture}
The unified architecture consists of:
\begin{enumerate}
    \item \textbf{Data Pipeline Scripts:} Python-based extraction, transformation, and verifiable provenance assignment.
    \item \textbf{Model Engine:} \texttt{scikit-learn} and \texttt{xgboost} implementations ensuring robust predictive capabilities via Grid Search CV.
    \item \textbf{React + Vite Dashboard:} A responsive, component-driven user interface featuring interactive spatial validations via \texttt{react-leaflet} and metric visualizations for lending officers.
\end{enumerate}

\section{Conclusion}
TerraTrust successfully demonstrates the integration of complex geospatial and environmental APIs into a unified, actionable financial metric. By strictly avoiding synthetic geographic parameters, the system proves its viability as an academically rigorous and verifiable tool for modern rural banking infrastructures.

\end{document}
